\documentclass[11pt,a4paper]{article}
\usepackage[margin=2.2cm]{geometry}
\usepackage[T1]{fontenc}
\usepackage[utf8]{inputenc}
\usepackage{charter,microtype,amsmath,amssymb,enumitem,xurl}
\usepackage[colorlinks=true,urlcolor=blue,pdftitle={Finite-sample Bayesian learning near singularities},pdfauthor={Francis J. Headley}]{hyperref}
\usepackage{iftex}
\ifPDFTeX\input{glyphtounicode}\pdfgentounicode=1\fi
\setlength{\parindent}{0pt}\setlength{\parskip}{6pt}
\setlist[itemize]{leftmargin=16pt,topsep=2pt,itemsep=2pt,parsep=0pt}
\begin{document}
\begin{center}
{\Large\bfseries Finite-sample Bayesian learning\\near singularities}\\[7pt]
Francis J. Headley\quad |\quad Research in progress\quad |\quad October 2026
\end{center}
\vspace{3pt}
\textbf{Question.} When does a Bayesian model begin to resolve a weak feature, and what does a learning-coefficient measurement tell us about that change?

I study this in reduced-rank regression, where an entire layer can be integrated out exactly. The aim is to understand the finite-sample behaviour of a tractable singular model before using similar diagnostics to interpret larger networks.

\textbf{Why this might matter for AI safety.} Complexity measurements are sometimes used to study learned structure. Their interpretation depends on what they can resolve. A weak feature can be visible to one posterior but suppressed at another sampling temperature; a change in a complexity statistic can also differ from a change in posterior mass in a chosen region. Controlled examples help separate these effects. The project has not yet tested them in a trained language model or demonstrated a safety intervention.

\textbf{Model and calculation.} Let $y=ABx+\varepsilon$, with $A\in\mathbb R^{N\times H}$, $B\in\mathbb R^{H\times M}$, independent standard Gaussian parameter priors, isotropic inputs and unit Gaussian observation noise. For true matrix $C$, the population evidence is
\[
Z_n(C)=\mathbb E_{A,B}\exp\!\left(-\frac n2\|AB-C\|_F^2\right).
\]
Integrating out $B$ gives
\[
Z_n(C)=\mathbb E_A\!\left[\det(I_H+nA^TA)^{-M/2}
\exp\!\left(-\frac n2\operatorname{tr}\!\left[CC^T(I_N+nAA^T)^{-1}\right]\right)\right].
\]
For rank-one truth $C=\delta uv^T$, the natural signal scale is $s=n\delta^2/2$. The project derives a limiting profile from a two-subspace average and a written uniform error bound near zero truth when the dominant stratum is unique. At a tie, a separate rank-one argument gives the correction
\[
\frac{Z_n(C)}{Z_n(0)}=
\Phi(s)\frac{\log[n/(1+s)]+O(1)}{\log n+O(1)},\qquad 0\le s\le n/2.
\]
These are statements about evidence ratios. Posterior-mass crossovers are checked separately using explicitly defined windows.

\textbf{What the checks show.} For the amplitude window $P(|\alpha-\delta|<\delta/2)=1/2$ in the $(N,H,M)=(2,1,2)$ model, the measured crossing is $s=1.75763,\ 1.64693,\ 1.57366$ at $n=10^4,10^8,10^{32}$. Grid refinement leaves these values stable. The slow drift is consistent with a logarithmic correction. This marker measures the projection along the target, not full function error.

\newpage
\textbf{Interpretation.} A fixed-truth asymptotic expansion can hide constants that depend strongly on the distance to a singularity. This makes it unsafe to infer every finite-sample transition by dropping those constants from a comparison involving $\lambda\log n$. The window definition and the limiting regime must be specified.

For the unlocalised tempered posterior in this model, inverse temperature $\beta$ replaces $n$ by $n\beta$. This gives a tractable example of a measurement whose resolution depends on temperature. Connecting it to a practical local learning-coefficient estimator additionally requires accounting for localisation, baseline subtraction and sampling error.

\textbf{Evidence and limitations.}
\begin{itemize}
\item A fresh independent scalar/vector Gaussian integral and fixed-truth derivative check agreed with the radial evaluator to about $4\times10^{-10}$ at the selected points.
\item Algebraic checks found no discrepancies in 8,000 dimension triples for the zero/rank-one coefficient and multiplicity calculations. These are consistency checks against a second implementation, not a substitute for proof.
\item The review found a missing transpose in the T1 sampled-data covariance generator. It is corrected and covered by a regression check. The old sampled-data T1 outputs are retained as historical results and are excluded from the current evidence pending a rerun.
\item The general tangent-cone proposal and the higher-rank tie upper bound remain open. Not every proof step has been independently checked.
\end{itemize}

\textbf{Prior work.} Near-singularity scaling and the one-unit model have substantial prior art in Watanabe and Amari (2002). Chen and Murfet (2025) discuss the effect of inverse temperature on the resolution of learning-coefficient estimates. The intended contribution here is a more explicit finite-sample calculation, uniform control near a singular truth, and careful tests of particular transition markers. An exhaustive priority review has not been completed.

\textbf{Next experiment.} Use a small trained model with an independently measurable weak feature. Sweep the estimator temperature and compare its trace with direct measurements of that feature, controlling for localisation and sampler convergence. This would test whether the resolution effect survives outside the exactly tractable posterior.

\textbf{Research materials.} The accompanying repository contains the written proofs of the uniform bound and the tie correction, the exact evaluator in Python, and a deterministic check script, \path{review_checks/check_calculations.py}, that reproduces the numbers above.

\textbf{AI assistance.} The project uses AI assistance for derivation checks, numerics and code review. These checks are not an independent human peer review; the proofs and code are included so that the claims can be inspected.

\vspace{3pt}
{\small
\textbf{References}\par
S. Watanabe and S. Amari (2002), \href{https://proceedings.neurips.cc/paper_files/paper/2002/file/c2ba1bc54b239208cb37b901c0d3b363-Paper.pdf}{The Effect of Singularities in a Learning Machine when the True Parameters Do Not Lie on such Singularities}, NeurIPS 15.\par
Z. Chen and D. Murfet (2025), \href{https://arxiv.org/abs/2504.18048}{Modes of Sequence Models and Learning Coefficients}, arXiv:2504.18048.
}
\end{document}
